\documentclass[11pt]{article}

\usepackage[margin=1in]{geometry}
\usepackage[T1]{fontenc}
\usepackage{lmodern}
\usepackage{microtype}
\usepackage{amsmath,amssymb,amsthm,mathtools}
\usepackage{booktabs}
\usepackage{enumitem}
\usepackage{etoolbox}
\usepackage{xcolor}
\usepackage[hidelinks]{hyperref}
\usepackage[nameinlink,noabbrev]{cleveref}

\definecolor{darkblue}{RGB}{25,61,98}
\hypersetup{
  colorlinks=true,
  linkcolor=darkblue,
  citecolor=darkblue,
  urlcolor=darkblue,
  pdftitle={The Theoretical Model in Aiyagari (1994)},
  pdfauthor={Summary prepared from Aiyagari (1994) and Aiyagari (1993a)}
}

\newtheoremstyle{propertystyle}
  {0.8em}{0.8em}{\normalfont}{}% space above/below, body font, indent
  {\bfseries\color{darkblue}}{.}{0.5em}{}
\theoremstyle{propertystyle}
\newtheorem{property}{Property}
\AtEndEnvironment{property}{\par}

\newcommand{\E}{\mathbb{E}}
\newcommand{\R}{\mathbb{R}}
\newcommand{\1}{\mathbf{1}}
\newcommand{\dd}{\,\mathrm{d}}
\newcommand{\supp}{\operatorname{supp}}
\newcommand{\as}{\text{a.s.}}

\title{\textbf{The Theoretical Model in Aiyagari (1994)}\\[0.4em]
\large Uninsured Idiosyncratic Risk, Borrowing Constraints, and Aggregate Saving}
\author{A notation-consistent summary of Aiyagari (1994)\\
and the propositions in Aiyagari (1993a)}
\date{}

\begin{document}
\maketitle

\begin{abstract}
This note reconstructs the economic model in Aiyagari's ``Uninsured
Idiosyncratic Risk and Aggregate Saving.''  It separates primitives and
assumptions, the definition of stationary recursive competitive equilibrium,
and the model's qualitative properties.  The properties that the published
article attributes to Aiyagari (1993a) are restated in the notation of the
1994 article.  No proofs are included.  Journal page references refer to
Aiyagari (1994); working-paper page references refer to Aiyagari (1993a).
\end{abstract}

\tableofcontents

\section{Definitions and assumptions}
\label{sec:definitions}

\subsection{Environment and markets}

Time is discrete, $t=0,1,2,\ldots$.  There is a unit mass of ex ante identical,
infinitely lived households.  Markets are incomplete: households can trade
only a one-period risk-free asset, which is physical capital in the production
economy.  There are no state-contingent claims and hence no direct insurance
against individual labor-endowment shocks.  A law of large numbers eliminates
aggregate uncertainty, so all aggregate quantities are constant in a
stochastic steady state even though individual quantities fluctuate
\cite[pp.~664--665]{aiyagari1994}.

Each household receives an idiosyncratic labor endowment $l_t$.  In the
expository model, $\{l_t\}$ is i.i.d. with distribution $F$ and bounded support
\[
  \supp(F)=[l_{\min},l_{\max}], \qquad l_{\min}>0, \qquad \E[l_t]=1.
\]
Shocks are independent across households.  Equivalently, a household supplies
one unit of time with idiosyncratic productivity $l_t$.  Aggregate effective
labor is therefore one.  The quantitative specification replaces the i.i.d.
process with a finite-state Markov chain; see \cref{subsec:quant}.

\subsection{Preferences and the household budget}

The representative household type has preferences
\begin{equation}
  \E_0\sum_{t=0}^{\infty}\beta^t U(c_t),
  \qquad 0<\beta<1,
  \label{eq:utility}
\end{equation}
where $c_t\geq 0$ is consumption.  The paper defines the rate of time
preference as
\begin{equation}
  \lambda\equiv\frac{1-\beta}{\beta}>0,
  \qquad \beta(1+\lambda)=1.
  \label{eq:lambda}
\end{equation}
For the recursive results invoked in the exposition, $U$ is taken to be
continuously differentiable, strictly increasing, and strictly concave; the
working-paper appendix also states boundedness as a sufficient condition for
its Bellman formulation \cite[p.~12, n.~17]{aiyagari1993a}.  Some results below
add a bound on relative risk aversion.

Let $a_t$ be beginning-of-period net asset holdings, $r$ the constant net real
return, and $w$ the wage per efficiency unit.  The budget constraint is
\begin{equation}
  c_t+a_{t+1}=w l_t+(1+r)a_t,
  \qquad c_t\geq0, \qquad a_t\geq-b \quad \as,
  \label{eq:budget}
\end{equation}
where $b\geq0$ is the maximum amount of debt.  Thus, a larger $b$ permits more
borrowing; $b=0$ prohibits borrowing.

For $r>0$, nonnegative consumption and present-value budget balance imply the
natural debt limit $a_t\geq-wl_{\min}/r$.  It is therefore without loss to write
\begin{align}
  a_t&\geq-\phi, \label{eq:effective-limit}\\
  \phi(b,w,r)&\equiv
  \begin{cases}
    \min\{b,wl_{\min}/r\},&r>0,\\
    b,&r\leq0.
  \end{cases}
  \label{eq:phi}
\end{align}
Aiyagari calls \cref{eq:phi} the \emph{fixed borrowing-limit} specification.
The alternative \emph{present-value borrowing limit} sets
\begin{equation}
  \phi^{\mathrm{pv}}(w,r)=\frac{w l_{\min}}{r}, \qquad r>0.
  \label{eq:pv-limit}
\end{equation}

It is useful to shift assets by the effective debt limit and define total
resources:
\begin{align}
  \widehat a_t&\equiv a_t+\phi\geq0, \label{eq:ahat}\\
  z_t&\equiv w l_t+(1+r)\widehat a_t-r\phi. \label{eq:z}
\end{align}
Then
\begin{align}
  c_t+\widehat a_{t+1}&=z_t,
  &c_t&\geq0,&\widehat a_{t+1}&\geq0, \label{eq:shifted-budget}\\
  z_{t+1}&=w l_{t+1}+(1+r)\widehat a_{t+1}-r\phi. \label{eq:z-transition}
\end{align}
The minimum feasible resource level is
\begin{equation}
  z_{\min}\equiv w l_{\min}-r\phi\geq0.
  \label{eq:zmin}
\end{equation}

\subsection{The income-fluctuation problem}

For the i.i.d. shock case, the value of total resources $z$ is the unique
solution (under the stated regularity conditions) to
\begin{equation}
  V(z;b,w,r)=
  \max_{0\leq\widehat a'\leq z}
  \left\{
    U(z-\widehat a')+
    \beta\int V\!\left(wl'+(1+r)\widehat a'-r\phi;b,w,r\right)
      \dd F(l')
  \right\}.
  \label{eq:bellman}
\end{equation}
The single-valued policy for shifted next-period assets is
\begin{equation}
  \widehat a'=A(z;b,w,r),
  \qquad c(z;b,w,r)=z-A(z;b,w,r).
  \label{eq:policies}
\end{equation}
Together with \cref{eq:z-transition}, this policy induces a Markov process for
$z$.  Let $\mu_z(\cdot;b,w,r)$ denote its stationary distribution when it
exists.  Long-run mean \emph{net} asset holdings are
\begin{equation}
  E_{a,w}(b,w,r)
  \equiv
  \int A(z;b,w,r)\,\mu_z(\dd z;b,w,r)-\phi(b,w,r).
  \label{eq:Eaw}
\end{equation}
When the equilibrium wage is substituted, write
\begin{equation}
  E_a(r)\equiv E_{a,w}\bigl(b,w(r),r\bigr).
  \label{eq:Ea}
\end{equation}
This is household capital supply in the production economy.

\subsection{Firms and aggregate technology}

Competitive firms operate a constant-returns neoclassical technology
$F(K,L)$.  With aggregate labor normalized to $L=1$, define per-capita output
as $f(k,1)$, where $k=K/L$.  Capital depreciates at rate $\delta$.  Profit
maximization gives
\begin{equation}
  r=f_1(k,1)-\delta,
  \qquad w=f_2(k,1).
  \label{eq:firm-foc}
\end{equation}
The capital-demand schedule $K(r)$ is the solution to the first condition in
\cref{eq:firm-foc}; the implied wage is denoted $w(r)$.

In the quantitative model,
\begin{equation}
  f(k,1)=k^{\alpha},
  \qquad
  K(r)=\left(\frac{\alpha}{r+\delta}\right)^{1/(1-\alpha)},
  \qquad
  w(r)=(1-\alpha)K(r)^{\alpha}.
  \label{eq:cobb-douglas}
\end{equation}

\subsection{Quantitative specialization}
\label{subsec:quant}

Section IV of the paper takes one period to be one year and uses
\begin{equation}
  \beta=0.96,
  \qquad \alpha=0.36,
  \qquad \delta=0.08,
  \qquad b=0.
  \label{eq:calibration}
\end{equation}
Utility is constant relative risk aversion (CRRA):
\begin{equation}
  U(c)=
  \begin{cases}
    \dfrac{c^{1-\gamma}-1}{1-\gamma},&\gamma\neq1,\\[0.7em]
    \log c,&\gamma=1,
  \end{cases}
  \qquad \gamma\in\{1,3,5\}.
  \label{eq:crra}
\end{equation}
Here $\gamma=-cU''(c)/U'(c)$ is relative risk aversion.  The paper denotes this
coefficient by $\mu$; $\gamma$ is used here to avoid confusing it with a
stationary distribution.

The labor endowment is approximated by a seven-state Markov chain that matches
\begin{equation}
  \log l_t=\rho\log l_{t-1}
  +\sigma\sqrt{1-\rho^2}\,\varepsilon_t,
  \qquad \varepsilon_t\sim N(0,1),
  \label{eq:labor-ar1}
\end{equation}
with $\sigma\in\{0.2,0.4\}$ and
$\rho\in\{0,0.3,0.6,0.9\}$.  The discrete endowment levels are rescaled so the
stationary mean of $l$ is one \cite[pp.~675--676]{aiyagari1994}.  With serial
correlation, the household state is $(a,l)$ (or equivalently $(z,l)$), the
value function and policy depend on the current $l$, and the expectation in
\cref{eq:bellman} uses the transition matrix $P(l'\mid l)$.

\section{Stationary recursive competitive equilibrium}
\label{sec:equilibrium}

Fix the primitives $\{\beta,U,F,\delta,b,P\}$.  A stationary recursive
competitive equilibrium is
\[
  \bigl(r,w,V,g_c,g_a,\mu,k\bigr),
\]
where $g_c(a,l)$ and $g_a(a,l)$ are the consumption and next-asset policies and
$\mu$ is a stationary cross-sectional probability distribution over household
states, such that the following conditions hold.

\paragraph{1. Household optimization.}
Given $(r,w)$, $V$ solves the household Bellman problem and the policies attain
its maximum.  For the general Markov endowment process this problem is
\begin{equation}
  V(a,l)=
  \max_{a'\geq-\phi}
  \left\{
    U\bigl(wl+(1+r)a-a'\bigr)
    +\beta\sum_{l'}P(l'\mid l)V(a',l')
  \right\},
  \label{eq:markov-bellman}
\end{equation}
where the choice must also make consumption nonnegative.  In the original
asset notation, the resulting policies satisfy
\begin{equation}
  g_c(a,l)+g_a(a,l)=wl+(1+r)a,
  \qquad g_c(a,l)\geq0,\qquad g_a(a,l)\geq-\phi.
  \label{eq:hh-opt}
\end{equation}

\paragraph{2. Stationarity of the cross-sectional distribution.}
Let $Q(\cdot\mid a,l)$ be the transition kernel generated by $g_a$ and
$P(l'\mid l)$.  The distribution $\mu$ is invariant:
\begin{equation}
  \mu(B)=\int Q(B\mid a,l)\,\mu(\dd a,\dd l)
  \quad\text{for every measurable set }B.
  \label{eq:invariant}
\end{equation}
Thus the cross section is constant although individual histories are not.

\paragraph{3. Firm optimization.}
Factor prices equal marginal products:
\begin{equation}
  r=f_1(k,1)-\delta,
  \qquad w=f_2(k,1).
  \label{eq:firm-equilibrium}
\end{equation}

\paragraph{4. Market clearing and aggregation.}
Effective labor and assets clear their markets:
\begin{equation}
  \int l\,\mu(\dd a,\dd l)=1,
  \qquad
  k=\int a\,\mu(\dd a,\dd l)=E_a(r)=K(r).
  \label{eq:market-clearing}
\end{equation}
The last equality is the paper's graphical steady-state condition: capital
supplied by households equals capital demanded by firms.  Aggregate goods
feasibility in the no-growth stationary equilibrium is
\begin{equation}
  C+\delta k=f(k,1),
  \qquad
  C\equiv\int g_c(a,l)\,\mu(\dd a,\dd l).
  \label{eq:resource-constraint}
\end{equation}
Gross investment is $\delta k$, net aggregate saving is zero, and the gross
saving rate reported by Aiyagari is
\begin{equation}
  s=\frac{\delta k}{f(k,1)}.
  \label{eq:saving-rate}
\end{equation}

\paragraph{Full-insurance benchmark.}
If idiosyncratic earnings risk is absent or perfectly insured, all households
are identical.  The steady-state return equals the time-preference rate and
capital is at the modified golden-rule level:
\begin{equation}
  r^{\mathrm{FI}}=\lambda,
  \qquad f_1(k^{\mathrm{FI}},1)-\delta=\lambda.
  \label{eq:full-insurance}
\end{equation}

\paragraph{Pure-exchange government-debt interpretation.}
Let per-capita government debt be $d$ and let the equal lump-sum tax financing
interest be $\tau=rd$.  The household budget becomes
\begin{equation}
  c_t+a_{t+1}=wl_t-rd+(1+r)a_t,
  \label{eq:debt-budget}
\end{equation}
and asset-market clearing becomes
\begin{equation}
  \int a\,\dd\mu=d.
  \label{eq:debt-clearing}
\end{equation}

\paragraph{Monetary interpretation.}
Let nominal money per capita be one, $m_t$ household nominal balances, and $p$
the price level.  Interpreting
\begin{equation}
  a_t=\frac{m_t-1}{p},
  \qquad b=\frac1p,
  \label{eq:money-map}
\end{equation}
makes $a_t\geq-b$ equivalent to $m_t\geq0$.  The stationary monetary-market
condition is $\int a\,\dd\mu=0$, equivalently $\int m\,\dd\mu=1$.

\section{Properties}
\label{sec:properties}

This section contains statements only; no proofs are reproduced.  ``Proposition
$j$'' in a source note refers to the numbered appendix proposition in Aiyagari
(1993a), translated into the definitions of \cref{sec:definitions}.

\subsection{Household problem and policy functions}

\begin{property}[Natural debt limit and no-Ponzi equivalence]
If $r>0$, then
\[
  a_t\geq-\frac{wl_{\min}}r\ \as\ \text{for every }t
  \quad\Longleftrightarrow\quad
  \lim_{t\to\infty}\frac{a_{t+1}}{(1+r)^t}\geq0\ \as.
\]
Consequently, nonnegative consumption and present-value budget balance imply
the natural debt limit.  \emph{Source: Aiyagari (1993a), Proposition 1,
p.~37; Aiyagari (1994), p.~666.}
\end{property}

\begin{property}[Value and policy regularity]
Under the regularity conditions stated in \cref{sec:definitions}, the Bellman
problem has a unique value function.  The shifted asset policy $A(z;b,w,r)$ is
single-valued and continuous.  For $z>z_{\min}$, consumption is positive,
$V$ is continuously differentiable in $z$, and
\[
  V'(z)=U'(c(z))
  \geq
  \beta(1+r)\E\bigl[V'(z')\mid z\bigr],
\]
with equality when $A(z;b,w,r)>0$.  \emph{Source: Aiyagari (1993a),
Proposition 2, pp.~37--38; Aiyagari (1994), pp.~666--667.}
\end{property}

\begin{property}[Binding region at low resources]
Suppose either $U'(0)<\infty$ or $z_{\min}>0$.  There is a cutoff
$\widehat z>z_{\min}$ such that, for every $z\leq\widehat z$,
\[
  c(z)=z,
  \qquad A(z;b,w,r)=0,
  \qquad a'=-\phi.
\]
Thus the household exhausts its borrowing capacity.  If $U'(0)=\infty$ and
$z_{\min}=0$, the borrowing limit is never exactly exhausted.
\emph{Source: Aiyagari (1993a), Proposition 3, p.~38; Aiyagari (1994),
p.~667.}
\end{property}

\begin{property}[Monotonicity away from the binding region]
For $z>\widehat z$, consumption and shifted next-period assets are strictly
increasing in $z$.  The asset policy satisfies
\[
  0<\frac{\partial A(z;b,w,r)}{\partial z}<1
\]
where differentiable.  In the binding region, a transitory endowment
innovation passes one-for-one into consumption.  \emph{Source: Aiyagari
(1994), p.~667 and n.~17.}
\end{property}

\begin{property}[Effect of a mean-preserving spread]
If $U'$ is convex, a mean-preserving spread in the labor-endowment distribution
lowers the consumption policy and raises $A(z;b,w,r)$ at every resource level.
The effect on stationary mean assets $E_{a,w}$ is not signed in general because
the invariant distribution also changes.  \emph{Source: Aiyagari (1994),
p.~670, n.~23, citing Sibley (1975) and Miller (1976).}
\end{property}

\subsection{Boundedness, stationarity, and aggregate asset supply}

\begin{property}[Bounded resource dynamics]
Let $R=1+r$.  If $\beta R<1$ (equivalently $r<\lambda$), the labor shock has
bounded support, and relative risk aversion
\[
  -\frac{cU''(c)}{U'(c)}
\]
is bounded above for all sufficiently large $c$, then there exists a finite
$z^*$ such that
\[
  z_t\geq z^* \quad\Longrightarrow\quad z_{t+1}\leq z_t\ \as.
\]
\emph{Source: Aiyagari (1993a), Proposition 4, pp.~38--39; Aiyagari
(1994), p.~668, n.~18.}
\end{property}

\begin{property}[Unique stable invariant distribution]
Under the assumptions of the preceding property, the Markov process for $z$
has a unique invariant distribution.  It is stable: from any initial
distribution, iterating the transition kernel converges to that invariant
distribution.  The invariant distribution varies continuously with
$(b,w,r)$.  \emph{Source: Aiyagari (1993a), Proposition 5, pp.~39--40;
Aiyagari (1994), p.~668.}
\end{property}

\begin{property}[Continuity and possible nonmonotonicity of mean assets]
For $r<\lambda$, $E_{a,w}(b,w,r)$ is finite and continuous in $b$, $w$, and
$r$, but it need not be monotone in $r$.  After the equilibrium wage $w(r)$ is
substituted, $E_a(r)$ remains continuous and need not be monotone.
\emph{Source: Aiyagari (1994), pp.~668 and 671, nn.~19 and 24.}
\end{property}

\begin{property}[Explosion at the time-preference rate]
Mean desired assets obey
\[
  \lim_{r\uparrow\lambda}E_{a,w}(b,w,r)=+\infty.
\]
For $r\geq\lambda$, the infinitely lived household accumulates without bound,
so finite stationary mean asset supply does not exist.  Hence every stationary
general equilibrium with finite capital satisfies $r<\lambda$.
\emph{Source: Aiyagari (1994), pp.~668--669.}
\end{property}

\begin{property}[Uncertainty raises desired assets near the benchmark]
For $r<\lambda$ but not too far below $\lambda$, mean assets under uninsured
idiosyncratic risk exceed mean assets under certain earnings.  In the certainty
case, every household holds the lowest permissible level $-\phi$ for
$r<\lambda$; under incomplete markets, the stationary cross section also
contains high-resource households that save above the limit.  This comparison
does not require convex marginal utility.  If $r$ is sufficiently low,
everyone may remain constrained and the two asset levels coincide.
\emph{Source: Aiyagari (1994), pp.~669--670 and n.~23.}
\end{property}

\subsection{General-equilibrium comparisons and existence}

\begin{property}[Factor-price and capital signs relative to full insurance]
Relative to the full-insurance economy,
\[
  r^{\mathrm{IM}}<r^{\mathrm{FI}}=\lambda,
  \qquad
  k^{\mathrm{IM}}>k^{\mathrm{FI}}.
\]
Thus uninsured idiosyncratic risk combined with limited borrowing produces
precautionary capital accumulation above the modified golden-rule stock.
\emph{Source: Aiyagari (1994), pp.~665 and 671.}
\end{property}

\begin{property}[Saving-rate sign]
If $k/f(k,1)$ is increasing in $k$ (strictly increasing when the capital share
is below one), then
\[
  \frac{\delta k^{\mathrm{IM}}}{f(k^{\mathrm{IM}},1)}
  >
  \frac{\delta k^{\mathrm{FI}}}{f(k^{\mathrm{FI}},1)}.
\]
The gross saving rate is therefore higher under incomplete markets.  Net
aggregate saving is zero in either no-growth steady state.
\emph{Source: Aiyagari (1994), p.~671 and n.~27.}
\end{property}

\begin{property}[Equilibrium need not be unique]
The household capital-supply curve $E_a(r)$ need not be monotone.  Therefore,
the market-clearing equation $K(r)=E_a(r)$ can have multiple solutions; the
paper does not claim uniqueness of the stationary general equilibrium.
\emph{Source: Aiyagari (1994), p.~671, nn.~24--25.}
\end{property}

\begin{property}[Existence under alternative debt limits]
With the present-value limit $\phi^{\mathrm{pv}}=wl_{\min}/r$ and
$l_{\min}>0$,
\[
  \lim_{r\downarrow0}E_{a,w}(\phi^{\mathrm{pv}},w,r)=-\infty.
\]
Together with divergence to $+\infty$ as $r\uparrow\lambda$, this ensures at
least one steady state with a positive interest rate.  With a fixed borrowing
limit, a positive-rate steady state need not exist; a steady state with a
negative interest rate does exist in the case described by the paper.
If $l_{\min}=0$, the present-value limit is equivalent to no borrowing.
\emph{Source: Aiyagari (1994), p.~673 and n.~30.}
\end{property}

\begin{property}[General equilibrium disciplines precautionary saving]
Because $E_a(r)$ becomes arbitrarily sensitive as $r$ approaches $\lambda$
from below, partial equilibrium can generate arbitrarily large excess asset
demand by fixing $r$ sufficiently close to $\lambda$.  In general equilibrium,
$r$ adjusts through $K(r)=E_a(r)$, limiting this mechanism.  The quantitative
allocation of the effect between a lower return and a higher capital stock
depends on the elasticity of $K(r)$.  \emph{Source: Aiyagari (1994), p.~672
and n.~28.}
\end{property}

\subsection{Borrowing limits and risk}

\begin{property}[More borrowing reduces precautionary capital]
Under a fixed debt limit, when $r=-1$ assets equal $-b$.  When $r=0$, mean net
assets fall one-for-one with $b$.  More generally, increasing $b$ shifts the
household asset-supply schedule to the left over the range in which the fixed
limit binds.  It lowers equilibrium capital and raises the equilibrium return
toward $\lambda$.  Hence allowing borrowing makes the capital and saving-rate
increases caused by uninsured risk smaller than under $b=0$.
\emph{Source: Aiyagari (1994), p.~672.}
\end{property}

\begin{property}[No effect when the fixed limit is slack]
Changes in $b$ affect asset supply only where $b<wl_{\min}/r$.  If the
equilibrium lies in a region where the natural limit is tighter, a marginal
change in $b$ has no effect on household behavior or the steady state.
\emph{Source: Aiyagari (1994), p.~672, n.~29.}
\end{property}

\subsection{Growth, government debt, and money}

\begin{property}[Balanced-growth extension]
Suppose labor-augmenting productivity, the wage, and the borrowing limit grow
at gross rate $1+g$, and utility is isoelastic with elasticity of marginal
utility $\gamma$.  Then the full-insurance balanced-growth return is
\[
  r^{\mathrm{FI}}_g=(1+\lambda)(1+g)^{\gamma}-1,
\]
whereas the incomplete-markets return is strictly below this value.  Earnings,
resources, assets, consumption, and income normalized by their per-capita
levels have stationary distributions.  \emph{Source: Aiyagari (1994), p.~671,
n.~26.}
\end{property}

\begin{property}[Debt neutrality depends on the borrowing constraint]
In the pure-exchange government-debt model, debt neutrality fails under the
fixed limit \cref{eq:phi}.  If instead the limit adjusts for tax obligations,
\[
  a_t\geq-\frac{wl_{\min}-rd}{r},
\]
then debt neutrality holds: the equilibrium return, consumption distribution,
and distribution of assets net of government debt are invariant to $d$.
\emph{Source: Aiyagari (1994), pp.~673--674.}
\end{property}

\begin{property}[Upper bound in the monetary interpretation]
Under the present-value debt limit, let $r^{\#}$ be the return satisfying zero
mean net assets.  It supports the price level
\[
  p^{\#}=\frac{r^{\#}}{wl_{\min}}
  \qquad\left(\text{equivalently }b^{\#}=\frac1{p^{\#}}\right).
\]
Lowering the price level cannot support a return above $r^{\#}$ because the
borrowing constraint is then slack on the relevant portion of asset supply.
Raising the price level above $p^{\#}$ tightens the constraint, raises mean
assets near $r^{\#}$, and lowers the equilibrium return.  In particular,
monetary equilibria cannot place $r$ arbitrarily close to $\lambda$.
\emph{Source: Aiyagari (1994), p.~674.}
\end{property}

\section{Notation crosswalk and source map}

\begin{center}
\small
\begin{tabular}{@{}lll@{}}
\toprule
Object & Notation here & Location in Aiyagari (1994) \\
\midrule
Discount factor / time-preference rate & $\beta$, $\lambda=(1-\beta)/\beta$ & p.~665 \\
Consumption / net assets / labor shock & $c_t$, $a_t$, $l_t$ & p.~665 \\
Maximum debt / effective debt limit & $b$, $\phi$ & pp.~665--666 \\
Shifted assets / total resources & $\widehat a_t$, $z_t$ & p.~666 \\
Shifted asset policy & $A(z;b,w,r)$ & p.~667 \\
Stationary mean net assets & $E_{a,w}$, $E_a(r)$ & pp.~668, 670--671 \\
Capital demand / wage schedule & $K(r)$, $w(r)$ & pp.~670--671 \\
Capital / depreciation / saving rate & $k$, $\delta$, $\delta k/f(k,1)$ & pp.~670--671 \\
CRRA coefficient & $\gamma$ (paper: $\mu$) & p.~675 \\
Labor persistence and dispersion & $\rho$, $\sigma$ & p.~675 \\
\bottomrule
\end{tabular}
\end{center}

\begin{thebibliography}{9}

\bibitem{aiyagari1994}
S.~Rao Aiyagari.
\newblock Uninsured idiosyncratic risk and aggregate saving.
\newblock \emph{The Quarterly Journal of Economics}, 109(3):659--684, 1994.

\bibitem{aiyagari1993a}
S.~Rao Aiyagari.
\newblock Uninsured idiosyncratic risk and aggregate saving.
\newblock Federal Reserve Bank of Minneapolis, Research Department Working
  Paper 502, revised December 1993 (first draft August 1992).

\end{thebibliography}

\end{document}
